% SOURCE_AUTHORITY: sga5_fr_workpass.tex, lines 9389--9436 (LF-normalized unit).
% SOURCE_SHA256: 791F4EFFC5E02832D5D77ED03518C8156D6F07E4C8238B03545DB93D883FBB28
\subsubsection*{2.2.5}
Para enunciar o seguinte corolário, vamos colocar, para todo esquema $X$,
\[
A^*(X)=\bigoplus_i H^{2i}(X,\mu_X^{\otimes i})
\]
e vamos usar sistematicamente esta notação no futuro.

\begin{statement}{Corolário 2.2.6}
O homomorfismo de imagem inversa
\[
p^*:A^*(S)=\bigoplus_i H^{2i}(S,\mu_S^{\otimes i})
\longrightarrow
\bigoplus_i H^{2i}(P,\mu_P^{\otimes i})=A^*(P)
\]
é um morfismo injetivo de anéis unitários, e os elementos $1,\xi,\ldots,\xi^r$ formam uma base do $A^*(S)$-módulo $A^*(P)$. Em particular, $A^*(P)$ é um $A^*(S)$-módulo livre do mesmo posto que $E$.
\end{statement}

\section*{3. CLASSES DE CHERN}

Sejam $S$ um esquema e $E$ uma $\mathcal O_S$-módulo localmente livre de posto constante $r$. Denotando por $\check E$ o dual de $E$, é colocado $P=P_S(\check E)$. Vimos (2.2.6) que, se $\xi$ denota a classe de $\mathcal O_P(1)$ em $H^2(P,\mu)$, o $A^*(S)$-módulo $A^*(P)$ é livre de base $1,\xi,\ldots,\xi^{r-1}$. Em particular, por razões de homogeneidade, é imediato que existe uma única relação entre a forma como a
\[
\xi^r+c_1\xi^{r-1}+\cdots+c_i\xi^{r-i}+\cdots+c_r=0,
\tag*{3.1.}
\]
com $c_i\in H^{2i}(S,\mu^{\otimes i})$ para tudo $i\in[0,r]$.

Seja Agora $E$ um $\mathcal O_S$-módulo localmente livre qualquer. Para todo inteiro $r$, o conjunto $S_r$ dos pontos de $S$ em que $E$ tem posto $r$ é um aberto. Assim obtém-se uma partição
\[
S=\coprod_r S_r
\]
de $S$ em abertos desajuntos; em particular, para todo $A_S$-módulo $M$ e todo inteiro $i$, o produto morfismo das imagens inversas
\[
H^i(S,M)\longrightarrow \prod_r H^i(S_r,M/S_r)
\]
É uma biéção.

As considerações anteriores permitem justificar a seguinte definição.

\begin{statement}{Definição 3.2}
Sejam $S$ um esquema e $E$ um $\mathcal O_S$-módulo localmente livre. Para todo inteiro $i\ge0$, chama-se $i$-ésimo classe de Chern de $E$, e denota-se por $c_i(E)$, para o elemento de $H^{2i}(S,\mu^{\otimes i})$ cuja restrição, para todo inteiro $r\ge0$, para o aberto $S_r$ de $S$ formado pelos pontos onde $E$ tem posto $r$ é o elemento $c_i$ definido pela relação 3.1. Se o posto de $E$ estiver acoplado, chama-se classe de Chern total de $E$ para o elemento
\[
c(E)=\sum_i c_i(E)
\]
de $A^*(S)$; em geral, designa-se assim o elemento correspondente de
\[
\widehat A^{\bullet}(S)=\prod_i A^i(S).
\]
\end{statement}
